\input{../../../headers/tpheaders_2.tex}

\prob{Modéliser un Système Linéaire Continu et Invariant} \\

\graphicspath{{../../../img/}}
\begin{center}
\def\svgwidth{\columnwidth}
\input{../../../img/triptyque.pdf_tex}
\end{center}

La démarche de l’ingénieur permet :
\begin{itemize}
 \item De vérifier les performances attendues d’un système, par évaluation de l’écart entre un cahier des charges et les réponses expérimentales (écart 1),
 \item De proposer et de valider des modèles d’un système à partir d’essais, par évaluation de l’écart entre les performances mesurées et les performances simulées (écart 2),
 \item De prévoir le comportement à partir de modélisations, par l’évaluation de l’écart entre les performances simulées et les performances attendues du cahier des charges (écart 3).
\end{itemize}

\documentsressource{Pour ce TP, vous aurez à votre disposition :
\begin{itemize}
 \item la \miseenoeuvre du système,
 \item le \href{https://github.com/Costadoat/Sciences-Ingenieur/raw/master/S02\%20Syst\%C3\%A8mes\%20Lin\%C3\%A9aires\%20Continus\%20Invariants/C01\%20Pr\%C3\%A9sentation\%20des\%20SLCI/02-C01.py}{script python} permettant de tracer la réponse indicielle à une fonction de transfert,
 \item les diverses informations sur le \urlsysteme.
\end{itemize}}
\newpage

\newpage

\section{Tracer la réponse à un échelon}

L'objectif de cette première partie est de tracer et d'analyser la réponse à un échelon pour un système.

~\

\question{En utilisant le logiciel d'acquisition, sans corde, appuyer sur le bouton de mise sous tension de la corde et tracer la \textbf{tension aux bornes du moteur} et la \textbf{vitesse du charriot} en fonction du \textbf{temps}.}

\question{Déterminer les performances du système dans la mesure du possible (précision, stabilité, rapidité,...).}

\question{A l'aide de vos connaissances et en vous aidant du cours, déterminer l'ordre et la classe de cette fonction de transfert.}

\question{A l'aide de vos connaissances et en vous aidant du cours, déterminer les valeurs numériques des caractéristiques de cette fonction de transfert.}


\section{Tracé simulé de la fonction de transfert déterminée}

\question{A l'aide du script python, tracer la fonction de transfert ainsi déterminée.}

\question{Comparer avec le tracé de la fonction précédente.}


\end{document}
